\chapter{Guide d'installation}
\label{annexe:installation}

SentinelOps s'installe de la même façon sur un PC, un VPS ou une instance cloud (AWS EC2,
Azure, Google Cloud, OVHcloud, Hetzner, DigitalOcean\ldots). Seuls la création de la machine
et le pare-feu dépendent du fournisseur.

\section*{Prérequis}
\begin{itemize}
  \item Linux (Ubuntu 22.04 ou 24.04 conseillé), 2 vCPU et 2 Go de RAM au minimum, 10 Go de
        disque ; sous Windows ou macOS : Docker Desktop ;
  \item Docker Engine avec le module Compose v2, Git et Python 3.8 ou plus.
\end{itemize}

\section*{Étapes sur un serveur Linux}
\begin{lstlisting}[language=bash,caption={Installation sur un serveur Linux},label={lst:install}]
# 1. Docker (script officiel) + droits pour l'utilisateur courant
curl -fsSL https://get.docker.com | sudo sh
sudo usermod -aG docker $USER && newgrp docker

# 2. Recuperer le projet
git clone https://github.com/Amin-BenZid/SentinelOps.git
cd SentinelOps

# 3. Demarrer
./sen run                 # acces par tunnel SSH (recommande)
./sen run --public        # ou acces direct sur le port 4000
\end{lstlisting}

\section*{Accès au dashboard}
\begin{itemize}
  \item \textbf{Tunnel SSH} : depuis le poste de l'opérateur,
        \code{ssh -L 4000:localhost:4000 utilisateur@IP_DU_SERVEUR}, puis ouvrir
        \url{http://localhost:4000}.
  \item \textbf{Port public} : autoriser le port TCP 4000 dans le pare-feu du fournisseur
        (groupe de sécurité AWS, NSG Azure, règle VPC Google Cloud, pare-feu du VPS)
        \textbf{uniquement depuis l'adresse IP de l'opérateur}, puis ouvrir
        \code{http://IP_DU_SERVEUR:4000}.
\end{itemize}

Le mot de passe initial s'affiche à la fin de \code{sen run} (ou avec \code{sen password}) ; il
doit être changé à la première connexion.

